%% Figure block, shared by ../paper/main.tex and ../submission/*.tex so the caption has
%% exactly one source.
\begin{figure}[tbp]
\centering
\begin{minipage}[t]{0.485\linewidth}
\centering
\figorbox{fig3_gapping.pdf}
\caption{\textbf{At matched cost, a deterministic gap decorrelates curvature updates
exponentially in $m$ while random thinning only does so linearly} (\Cref{prop:gap}).
Vertical axis: $n\,\E\lVert\Hb-H\rVert_F^2$ divided by its exact independent-data value
$(\tr\Sigma)^2+\lVert\Sigma\rVert_F^2$ --- a computed constant, not a fitted normalisation.
Pale bands are the two predictions; markers are measurements with the optimisation frozen at
$\ts$, so only estimation quality is compared. $\kappa_H=(1+\phi^2)/(1-\phi^2)$ here.}
\label{fig:gap}
\end{minipage}\hfill
\begin{minipage}[t]{0.485\linewidth}
\centering
\figorbox{fig4_cost.pdf}
\caption{\textbf{The statistically correct interval is the cheaper one.} Wall-clock
milliseconds for one pass over $N=200{,}000$ observations (vertical axis, log scale) against
dimension $d$ (horizontal, log scale), single-threaded, minimum over seven interleaved
repetitions. The i.i.d.\ plug-in score covariance costs $O(d^2)$ per \emph{observation}; the
blocked long-run covariance costs $O(d^2)$ per \emph{block}. Slopes are fitted on the four
largest $d$; at this $N$ they are inflated by the one-off warm-up (see \eqref{eq:cost} and
\Cref{tab:cost}, which also gives the asymptotic exponents), and the two smallest $d$ are
dominated by fixed per-block overhead rather than by either term of \eqref{eq:cost}. The
message of the figure is the vertical separation, which holds at every $d$.}
\label{fig:cost}
\end{minipage}
\end{figure}
